\chapter[Conclusão]{Conclusão}
\label{chap: conclusao}

Este trabalho investigou se um algoritmo genético é capaz de equilibrar competitivamente cinco arquétipos de um jogo de luta sem destruir as suas identidades funcionais. Para isso, construiu um modelo de combate em que cada característica dos personagens tem efeito mensurável, formulou uma função de aptidão que contém o que cada arquétipo é, mas não quem deve vencer quem, e comparou um algoritmo genético escalar, o NSGA-II e um híbrido dos dois sob um protocolo com modelos nulos, controles e validação fora do laço de otimização.

A resposta é afirmativa, com uma qualificação. O algoritmo escalar alcança o equilíbrio em todas as execuções, com elencos tão equilibrados quanto a simetria perfeita, e o equilíbrio replica fora das condições de treino. A identidade funcional sobrevive apenas em parte, mas a parte que sobrevive é atribuível ao termo de identidade: sem ele, a mesma busca equilibra o elenco sem ganho significativo de equilíbrio e produz personagens cuja identidade não se distingue da de elencos sorteados. Identidade e equilíbrio não se mostraram antagônicos na faixa medida. E a identidade que o algoritmo escalar entrega não é a máxima compatível com o equilíbrio: com o mesmo orçamento, o híbrido alcança a identidade do NSGA-II mantendo o equilíbrio do algoritmo escalar, e a concordância comportamental com os arquétipos quase dobra.

\section{Contribuições}

O trabalho deixa cinco contribuições.

\begin{enumerate}
    \item \textbf{Uma formulação não circular do balanceamento com preservação de identidade.} A separação entre premissa e resposta permite incluir a identidade na busca sem codificar o resultado que se quer observar. A identidade estrutural entra como penalidade, que os resultados mostram poder perder, e a identidade funcional é medida por réguas comportamentais que a busca não vê.
    \item \textbf{Evidência controlada sobre o efeito do termo de identidade.} O contrafactual sem o termo, com a mesma amostra e o mesmo orçamento, isola o efeito do método: a identidade preservada se deve ao termo, que não custou equilíbrio mensurável.
    \item \textbf{A identificação de uma falha de busca sob avaliação ruidosa.} A seleção escalar, ao somar um termo exato com um termo amostrado, termina num extremo do compromisso e não no ótimo da sua própria função. O compromisso que ela exibe é um limite superior do custo da identidade, e a multiobjetivização, na forma do híbrido, recupera a identidade perdida com o mesmo orçamento de gerações e sem custo de equilíbrio.
    \item \textbf{A localização da perda de identidade na política.} Os parâmetros de comportamento, que o equilíbrio quase não enxerga, são os menos preservados em todos os métodos. A preservação de parâmetros invisíveis à função de equilíbrio depende inteiramente da pressão de identidade.
    \item \textbf{Lições de método para a avaliação por simulação.} Métricas de identidade precisam de piso medido; uma régua de medição grosseira pode favorecer sistematicamente um braço experimental; e a resolução da medida precisa acompanhar a margem do fenômeno medido. O sistema desenvolvido registra, em cada artefato, a configuração e o código que o produziram, o que torna cada número deste trabalho reprodutível e rastreável.
\end{enumerate}

\section{Trabalhos futuros}

As limitações discutidas na \autoref{sec: limitacoes} apontam as direções seguintes.

\textbf{Testar o equilíbrio contra oponentes que se adaptam.} O equilíbrio medido é condicionado a uma política fixa. Coevoluir estratégias que procurem explorar o elenco evoluído \cite{chen2014mmorpg, livingstone2006coevolution} verificaria se o equilíbrio resiste a quem o ataca, e não apenas a sorteios.

\textbf{Dar à política o peso que a identidade exige.} A política é o componente de identidade menos preservado. Mutação mais ampla nos pesos comportamentais, inicialização da população em torno dos perfis canônicos ou um termo de identidade que pondere a política de acordo com a sua baixa visibilidade para o equilíbrio são alternativas diretas, cujo efeito o protocolo deste trabalho já permite medir.

\textbf{Políticas dependentes do estado.} Uma política que responda à distância, aos pontos de vida e ao estado do oponente permitiria representar planos, e não apenas misturas de posturas, e aproximaria os arquétipos modelados dos de jogos reais.

\textbf{Mecânicas que sustentem ciclos.} O roster canônico se revelou uma hierarquia no modelo. Acrescentar mecânicas de que as vantagens clássicas dependem, como janelas de punição e encadeamento de golpes, testaria se o ciclo autoral passa a ser realizável e se o algoritmo o preserva quando ele existe.

\textbf{Explorar a diversidade de elencos equilibrados.} Em vez de procurar um elenco, métodos de qualidade e diversidade \cite{mouret2015mapelites} poderiam iluminar quantas configurações distintas de elenco atingem o equilíbrio, usando o perfil comportamental como espaço de características.

\textbf{Calibrar o híbrido.} A repartição do orçamento entre as fases do híbrido foi escolhida por simplicidade. Uma varredura no orçamento completo diria qual repartição melhor equilibra identidade, equilíbrio e velocidade de convergência.
